\chapter{Comparison to Racket's contract system}
\label{chap:compare}

Previously, we described an implementation contract system in OCaml. Our implementation relies heavily on the preprocessor rewriting engine to wrap a value or function with its contract. Guarded values are rewritten into a function to provide blame labels and checked when all labels are provided. In this section we provided a few insight into the Racket's contract system, then compare it to our approach in OCaml. We will see that implementing contract system in OCaml is a challenge compared to Racket due to OCaml's restriction.

\section{Racket's contract system}

The Racket language provides two main functionalities that makes implementing contract system easier: hygenic macros, and chaperone. The contract system of Racket\footnote{\href{https://github.com/racket/racket/tree/master/racket/collects/racket/contract}{https://github.com/racket/racket/tree/master/racket/collects/racket/contract}} is implemented using (hygenic) macros to modify the sourcecode. Afterwards, guarded values are wrapped inside a chaperone. Racket macros allows implementors to manipulate the AST in a more sensible manner during compile time. It replaces the need for a preprocessor, and target specific. Chaperone on the other hand, provides a special run-time object manipulation that is useful for contract system implementation.

A chaperone is an \textit{invisble proxy} over a value. When the value is used, the chaperone intercepts and perform some predefined actions. In the case of software contract, the contract checking is the chaperone action. A chaperone cannot do anything, for such is allowed by impersonator. The chaperone action can throw an exception, which is used to signal a contract violation. Chaperone cannot be used to implement flat contracts, but it can be used to implement function contracts using \texttt{chaperone-procedure}. This function takes a function and a wrapper function, which is applied for every argument, and return value(s). The wrapper function body should follows the (higher-order and dependent) functional contract semantics.

Chaperone is also lazy. A chaperone action on a vector or structure only runs when accessing the element (vector) or member (structure). This can be thought of as an optimization to lazily check until the value is actually needed.

\section{Comparison}

OCaml preprocessor and Racket macros are both used for the sole purpose of rewriting the AST to add contract checking functionalities to the sourcecode, eventhough OCaml provides a more generic system (preprocessor can modify any part of the AST). In both implementations, the guarded value must be wrapped. Racket's chaperone has support for function wrappers, lifting heavy and complex logic into a single function invocation (by \texttt{chaperone-procedure} and its family). OCaml does not have such functionalities, nor can it be implemented on top of the language (judging by Racket's choice of implementation\footnote{Chaperone is implemented as part of the language of Racket}). Therefore, functions are modified to wrapped each arguments and the body (return value).

Through our observation, we identify two main aspects that our implementation is greatly different from Racket's contract system: (i) Contract checking, and (ii) Blame assignment.

\subsection{Contract checking}

Our current implementation limits to values (as a whole) and functions. In Racket's contract system, contracts for vectors and structures is also implemented using chaperone to guard an element access (vector), or member access (structure). A chaperone enforces contract lazily, only when an element is accessed or on member access. Our implementation, however, perform contract checking for arrays and structs as an entity bringing in more runtime overhead.

Racket's chaperone can perform multi-shot contract checking, compared to our implementation where contract is checked only one time. For instance, a chaperone of vector guards its element across mutation; our contract system limits to one time contract checking. Vector, tuple and structure contracts in our implementation lacks this fine-grain control. This limitation is because OCaml's lack of chaperone-like mechanism.

One exception. Structures in OCaml, records, can be made into their functional versions using the preprocessor \textit{ppx\_fields\_conv}.\footnote{\href{https://github.com/janestreet/ppx\_fields\_conv}{https://github.com/janestreet/ppx\_fields\_conv}} Using the functions generated, we can implement our contract system to append contracts to these functions and provide contracts for structures. However, the original record dot notation is still unprotected.

\subsection{Blame assignment}

Blame assignment in Racket is handled by chaperone. By storing blame labels inside the ``proxy'' object, the protected values can be used without changing their invocation.

In our implementation, all protected values is converted into a function accepting three additional arguments for labels. Therefore, we need to search for all usages, and provide the correct blame labels. It becomes a hurdle when function signatures must also be modified.
